\subsection{Теорема о спектре самосопряженного оператора: \texorpdfstring{\(\sigma(A)\subset[m_-,m_+]\), \(r(A)=\|A\|\)}{sigma(A) subset [m-,m+], r(A)=||A||}}

\begin{definition}[Числа \(m_-\) и \(m_+\)]
	Пусть \(H\) --- комплексное гильбертово пространство, и
	\[
	A\in\mathcal{L}(H)
	\]
	--- самосопряженный оператор.
	
	Положим
	\[
	m_-=\inf_{\|x\|=1}(Ax,x),
	\qquad
	m_+=\sup_{\|x\|=1}(Ax,x).
	\]
	Так как \(A\) самосопряжен, то
	\[
	(Ax,x)\in\mathbb{R}.
	\]
	Кроме того, при \(\|x\|=1\)
	\[
	|(Ax,x)|\leqslant \|A\|.
	\]
	Поэтому числа \(m_-\) и \(m_+\) корректно определены и конечны.
\end{definition}

\begin{remark}[Полускалярное произведение]
	Если \(B\in\mathcal{L}(H)\) самосопряжен и
	\[
	(Bx,x)\geqslant0
	\qquad
	\forall x\in H,
	\]
	то формула
	\[
	\langle x,y\rangle_B=(Bx,y)
	\]
	задаёт полускалярное произведение.
	
	Для него верно неравенство Коши--Буняковского:
	\[
	|\langle x,y\rangle_B|^2
	\leqslant
	\langle x,x\rangle_B\langle y,y\rangle_B.
	\]
\end{remark}

\begin{remark}[Используемые факты]
	В доказательстве используем уже доказанные факты:
	
	\begin{itemize}
		\item спектр самосопряженного оператора вещественен:
		\[
		\sigma(A)\subset\mathbb{R};
		\]
		
		\item критерий принадлежности спектру:
		\[
		\lambda\in\sigma(A)
		\Longleftrightarrow
		\exists x_n:\|x_n\|=1,\quad (A-\lambda I)x_n\to0;
		\]
		
		\item основная теорема о спектральном радиусе:
		\[
		r(A)=\lim_{n\to\infty}\sqrt[n]{\|A^n\|}.
		\]
	\end{itemize}
\end{remark}

\begin{theorem}[Спектр самосопряженного оператора]
	Пусть \(H\) --- комплексное гильбертово пространство, и
	\[
	A\in\mathcal{L}(H)
	\]
	--- самосопряженный оператор. Тогда
	\[
	\sigma(A)\subset[m_-,m_+],
	\]
	причём
	\[
	m_-,m_+\in\sigma(A).
	\]
	Кроме того,
	\[
	r(A)=\max\{|m_-|,|m_+|\}=\|A\|.
	\]
\end{theorem}

\begin{proof}
	Докажем утверждение по смысловым шагам.
	
	\textbf{1. Спектр лежит внутри отрезка \([m_-,m_+]\).}
	
	Так как \(A\) самосопряжен, то по предыдущему билету
	\[
	\sigma(A)\subset\mathbb{R}.
	\]
	Поэтому достаточно доказать, что все вещественные \(\lambda\), лежащие вне отрезка
	\([m_-,m_+]\), являются регулярными.
	
	\begin{itemize}
		\item \textbf{Случай \(\lambda>m_+\).}
		
		Пусть
		\[
		\lambda>m_+.
		\]
		Для любого \(x\in H\), \(x\neq0\), имеем
		\[
		\|A_\lambda x\|\cdot\|x\|
		\geqslant
		|(A_\lambda x,x)|.
		\]
		Но
		\[
		(A_\lambda x,x)
		=
		((A-\lambda I)x,x)
		=
		(Ax,x)-\lambda\|x\|^2.
		\]
		По определению \(m_+\)
		\[
		(Ax,x)\leqslant m_+\|x\|^2.
		\]
		Так как \(\lambda>m_+\), получаем
		\[
		|(Ax,x)-\lambda\|x\|^2|
		=
		\lambda\|x\|^2-(Ax,x)
		\geqslant
		(\lambda-m_+)\|x\|^2.
		\]
		Следовательно,
		\[
		\|A_\lambda x\|\cdot\|x\|
		\geqslant
		(\lambda-m_+)\|x\|^2.
		\]
		Сокращая на \(\|x\|\), получаем
		\[
		\|A_\lambda x\|
		\geqslant
		(\lambda-m_+)\|x\|.
		\]
		Значит, \(A_\lambda\) ограничен снизу, поэтому по критерию из билета 11.3
		\[
		\lambda\in\rho(A).
		\]
		
		\item \textbf{Случай \(\lambda<m_-\).}
		
		Пусть
		\[
		\lambda<m_-.
		\]
		Тогда по определению \(m_-\)
		\[
		(Ax,x)\geqslant m_-\|x\|^2.
		\]
		Следовательно,
		\[
		|(Ax,x)-\lambda\|x\|^2|
		=
		(Ax,x)-\lambda\|x\|^2
		\geqslant
		(m_- - \lambda)\|x\|^2.
		\]
		Повторяя оценку через Коши--Буняковского, получаем
		\[
		\|A_\lambda x\|
		\geqslant
		(m_- - \lambda)\|x\|.
		\]
		Значит, \(A_\lambda\) ограничен снизу, и снова
		\[
		\lambda\in\rho(A).
		\]
	\end{itemize}
	
	Итак, все точки вне отрезка \([m_-,m_+]\) лежат в резольвентном множестве. Значит,
	\[
	\sigma(A)\subset[m_-,m_+].
	\]
	
	\textbf{2. Докажем, что \(m_+\in\sigma(A)\).}
	
	По определению \(m_+\) как супремума нужно построить последовательность почти собственных
	векторов для \(m_+\), а потом применить критерий из билета 11.3.
	
	\begin{itemize}
		\item \textbf{Выберем максимизирующую последовательность.}
		
		По определению супремума существует последовательность \(\{x_n\}\subset H\), такая что
		\[
		\|x_n\|=1
		\]
		и
		\[
		(Ax_n,x_n)\to m_+.
		\]
		
		\item \textbf{Введём неотрицательный самосопряженный оператор.}
		
		Рассмотрим оператор
		\[
		B_+=m_+I-A.
		\]
		Он самосопряжен. Кроме того, для любого \(x\in H\)
		\[
		(B_+x,x)
		=
		m_+\|x\|^2-(Ax,x)
		\geqslant0.
		\]
		Значит, формула
		\[
		\langle u,v\rangle_{B_+}=(B_+u,v)
		\]
		задаёт полускалярное произведение.
		
		\item \textbf{Применим Коши--Буняковского для полускалярного произведения.}
		
		Для этого полускалярного произведения имеем
		\[
		|\langle u,v\rangle_{B_+}|^2
		\leqslant
		\langle u,u\rangle_{B_+}\langle v,v\rangle_{B_+}.
		\]
		Берём
		\[
		u=x_n,
		\qquad
		v=B_+x_n.
		\]
		Тогда
		\[
		\|B_+x_n\|^4
		=
		|(B_+x_n,B_+x_n)|^2
		=
		|\langle x_n,B_+x_n\rangle_{B_+}|^2.
		\]
		По неравенству Коши--Буняковского
		\[
		\|B_+x_n\|^4
		\leqslant
		\langle x_n,x_n\rangle_{B_+}
		\langle B_+x_n,B_+x_n\rangle_{B_+}.
		\]
		
		\item \textbf{Покажем, что правая часть стремится к нулю.}
		
		Первый множитель равен
		\[
		\langle x_n,x_n\rangle_{B_+}
		=
		(B_+x_n,x_n)
		=
		m_+-(Ax_n,x_n)
		\to0.
		\]
		Второй множитель ограничен:
		\[
		\langle B_+x_n,B_+x_n\rangle_{B_+}
		=
		(B_+^2x_n,B_+x_n)
		\leqslant
		\|B_+^2x_n\|\cdot\|B_+x_n\|
		\leqslant
		\|B_+\|^3.
		\]
		Следовательно,
		\[
		\|B_+x_n\|^4\to0,
		\]
		то есть
		\[
		B_+x_n\to0.
		\]
		
		\item \textbf{Применим критерий принадлежности спектру.}
		
		Так как
		\[
		B_+x_n=(m_+I-A)x_n=-(A-m_+I)x_n,
		\]
		получаем
		\[
		(A-m_+I)x_n\to0.
		\]
		При этом
		\[
		\|x_n\|=1.
		\]
		По критерию из билета 11.3
		\[
		m_+\in\sigma(A).
		\]
	\end{itemize}
	
	\textbf{3. Докажем, что \(m_-\in\sigma(A)\).}
	
	Доказательство полностью аналогично предыдущему пункту, только вместо супремума берём
	инфимум.
	
	\begin{itemize}
		\item \textbf{Выберем минимизирующую последовательность.}
		
		По определению \(m_-\) существует последовательность \(\{y_n\}\subset H\), такая что
		\[
		\|y_n\|=1
		\]
		и
		\[
		(Ay_n,y_n)\to m_-.
		\]
		
		\item \textbf{Введём неотрицательный самосопряженный оператор.}
		
		Рассмотрим
		\[
		B_-=A-m_-I.
		\]
		Тогда \(B_-\) самосопряжен и
		\[
		(B_-x,x)
		=
		(Ax,x)-m_-\|x\|^2
		\geqslant0.
		\]
		
		\item \textbf{Повторим оценку из предыдущего пункта.}
		
		Для полускалярного произведения
		\[
		\langle u,v\rangle_{B_-}=(B_-u,v)
		\]
		аналогично получаем
		\[
		\|B_-y_n\|^4
		\leqslant
		\langle y_n,y_n\rangle_{B_-}
		\langle B_-y_n,B_-y_n\rangle_{B_-}.
		\]
		Здесь
		\[
		\langle y_n,y_n\rangle_{B_-}
		=
		(B_-y_n,y_n)
		=
		(Ay_n,y_n)-m_-
		\to0,
		\]
		а второй множитель ограничен сверху константой, не зависящей от \(n\).
		
		Значит,
		\[
		B_-y_n\to0.
		\]
		
		\item \textbf{Применим критерий принадлежности спектру.}
		
		Так как
		\[
		B_-y_n=(A-m_-I)y_n,
		\]
		получаем
		\[
		(A-m_-I)y_n\to0.
		\]
		При этом
		\[
		\|y_n\|=1.
		\]
		По критерию из билета 11.3
		\[
		m_-\in\sigma(A).
		\]
	\end{itemize}
	
	\textbf{4. Найдём спектральный радиус через \(m_-\) и \(m_+\).}
	
	Мы уже доказали, что
	\[
	\sigma(A)\subset[m_-,m_+],
	\]
	и при этом
	\[
	m_-,m_+\in\sigma(A).
	\]
	Поэтому
	\[
	r(A)
	=
	\sup_{\lambda\in\sigma(A)}|\lambda|
	=
	\max\{|m_-|,|m_+|\}.
	\]
	
	\textbf{5. Докажем равенство \(r(A)=\|A\|\).}
	
	Используем общую формулу спектрального радиуса:
	\[
	r(A)=\lim_{n\to\infty}\sqrt[n]{\|A^n\|}.
	\]
	Нужно показать, что для самосопряженного оператора этот предел равен \(\|A\|\).
	
	\begin{itemize}
		\item \textbf{Докажем равенство \(\|A^2\|=\|A\|^2\).}
		
		Оценка
		\[
		\|A^2\|\leqslant\|A\|^2
		\]
		следует из мультипликативности нормы оператора.
		
		Докажем обратную оценку. Для любого \(x\in H\), \(\|x\|=1\), имеем
		\[
		\|Ax\|^2
		=
		(Ax,Ax)
		=
		(A^2x,x).
		\]
		По неравенству Коши--Буняковского
		\[
		(A^2x,x)
		\leqslant
		\|A^2x\|\cdot\|x\|
		\leqslant
		\|A^2\|.
		\]
		Значит,
		\[
		\|Ax\|^2\leqslant\|A^2\|.
		\]
		Переходя к супремуму по всем \(\|x\|=1\), получаем
		\[
		\|A\|^2\leqslant\|A^2\|.
		\]
		Следовательно,
		\[
		\|A^2\|=\|A\|^2.
		\]
		
		\item \textbf{Перейдём к степеням \(2^k\).}
		
		Так как \(A\) самосопряжен, то \(A^2\), \(A^4\), \(\ldots\), \(A^{2^k}\) тоже самосопряжены.
		Применяя доказанное равенство последовательно, получаем
		\[
		\|A^{2^k}\|=\|A\|^{2^k}.
		\]
		
		\item \textbf{Применим формулу спектрального радиуса.}
		
		По общей формуле
		\[
		r(A)=\lim_{n\to\infty}\sqrt[n]{\|A^n\|}.
		\]
		Тогда по подпоследовательности \(n=2^k\)
		\[
		r(A)
		=
		\lim_{k\to\infty}\sqrt[2^k]{\|A^{2^k}\|}
		=
		\lim_{k\to\infty}\sqrt[2^k]{\|A\|^{2^k}}
		=
		\|A\|.
		\]
	\end{itemize}
	
	Итак,
	\[
	r(A)=\max\{|m_-|,|m_+|\}=\|A\|.
	\]
\end{proof}